\chapter{Scripts Docker}
\label{annexe:docker}

\section{Dockerfile appServer-1}

\begin{lstlisting}[language=bash, caption={Dockerfile pour le serveur d'application Flask}]
FROM python:3.11-slim

WORKDIR /opt

RUN pip install flask psycopg2-binary requests

COPY app.py .

EXPOSE 80

CMD ["python3", "app.py"]
\end{lstlisting}

\section{Application Flask (app.py)}

\begin{lstlisting}[language=python, caption={Serveur d'application Flask}]
from flask import Flask, request
import psycopg2, os

app = Flask(__name__)

DB_HOST = os.getenv("DB_HOST", "192.168.10.11")
DB_NAME = os.getenv("DB_NAME", "appdb")
DB_USER = os.getenv("DB_USER", "appuser")
DB_PASS = os.getenv("DB_PASS", "apppass")

@app.route("/")
def home():
    return "<h1>App Server</h1><p>Connected to PostgreSQL</p>"

@app.route("/health")
def health():
    return {"status": "ok", "service": "app-server"}

@app.route("/db-check")
def db_check():
    try:
        conn = psycopg2.connect(host=DB_HOST, dbname=DB_NAME,
                                user=DB_USER, password=DB_PASS)
        cur = conn.cursor()
        cur.execute("SELECT 1")
        return {"database": "connected", "result": cur.fetchone()[0]}
    except Exception as e:
        return {"database": "error", "message": str(e)}, 500

if __name__ == "__main__":
    app.run(host="0.0.0.0", port=80)
\end{lstlisting}

\section{Script dnsmasq (DHCP-Server)}

\begin{lstlisting}[language=bash, caption={Configuration dnsmasq pour le serveur DHCP}]
#!/bin/sh
ip addr add 192.168.20.2/24 dev eth0
ip link set eth0 up
ip route add default via 192.168.20.1

cat > /etc/dnsmasq.conf << EOF
interface=eth0
dhcp-range=192.168.20.100,192.168.20.200,255.255.255.0,12h
dhcp-option=3,192.168.20.1
dhcp-option=6,8.8.8.8
EOF

dnsmasq --no-daemon
\end{lstlisting}

\section{Configuration PostgreSQL}

\begin{lstlisting}[language=bash, caption={Initialisation de la base de donnees}]
# Demarrer PostgreSQL
su - postgres -c "pg_ctl -D /var/lib/postgresql/data start"

# Creer la base et l'utilisateur
su - postgres -c "psql -c \"CREATE DATABASE appdb;\""
su - postgres -c "psql -c \"CREATE USER appuser WITH PASSWORD 'apppass';\""
su - postgres -c "psql -c \"GRANT ALL PRIVILEGES ON DATABASE appdb TO appuser;\""

# Creer les tables
su - postgres -c "psql -d appdb -c \"
CREATE TABLE users (
    id SERIAL PRIMARY KEY,
    name VARCHAR(100),
    email VARCHAR(100)
);

CREATE TABLE requests (
    id SERIAL PRIMARY KEY,
    user_id INT REFERENCES users(id),
    description TEXT,
    created_at TIMESTAMP DEFAULT CURRENT_TIMESTAMP
);

INSERT INTO users (name, email) VALUES
    ('Alice', 'alice@example.com'),
    ('Bob', 'bob@example.com'),
    ('Charlie', 'charlie@example.com');
\""
\end{lstlisting}


\begin{thebibliography}{12}
    \bibitem{1} Fortinet. \textit{FortiOS 7.4 --- Administration Guide}. Fortinet Documentation, 2024. Disponible sur : \url{https://docs.fortinet.com}
    \bibitem{2} Fortinet. \textit{FortiGate --- High Availability Configuration Guide}. Fortinet Documentation, 2024.
    \bibitem{3} Fortinet. \textit{FortiGate --- IPSec VPN Configuration Guide}. Fortinet Documentation, 2024.
    \bibitem{4} Fortinet. \textit{FortiGate --- SSL VPN Configuration Guide}. Fortinet Documentation, 2024.
    \bibitem{5} Fortinet. \textit{FortiGate --- Security Profiles Guide}. Fortinet Documentation, 2024.
    \bibitem{6} GNS3 Documentation. \textit{GNS3 --- User Guide}. Disponible sur : \url{https://docs.gns3.com}
    \bibitem{7} Coursera. \textit{FortiGate Administrator --- Module 1 à 7}. Fortinet / Coursera, 2026.
    \bibitem{8} Coursera. \textit{Enterprise Firewall Administrator --- Module 1 à 5}. Fortinet / Coursera, 2026.
    \bibitem{9} Coursera. \textit{FortiAnalyzer Administrator --- Module 1 à 4}. Fortinet / Coursera, 2026.
    \bibitem{10} Moyen, J.-P. \textit{Les réseaux}. Eyrolles, 6e édition, 2012.
    \bibitem{11} Axeli. \textit{Site officiel --- Présentation de l'entreprise}. Disponible sur : \url{https://www.axeli.ma}
    \bibitem{12} CDIO Initiative. \textit{EMSI --- Moroccan School of Engineering Sciences}. Disponible sur : \url{https://www.cdio.org}
\end{thebibliography}
